\scanpage{261}
\mspage{7}
\margin{n° 324}
Cette fonctions est \uline{dérivable} en $x_0$, et sa dérivée est donnée par la formule :
\[
(7)\quad (f_1\cdot f_2)'(x_0)\cdot h=(f'_1(x_0)\cdot h)\cdot f_2(x_0)+f_1(x_0)\cdot(f'_2(x_0)\cdot h),\quad h\in E,
\]
formule que l'on écrit aussi plus brièvement :
\[
(7\text{ bis})\quad (f_1\cdot f_2)'(x_0)=f'_1(x_0)\cdot f_2(x_0)+f_1(x_0)\cdot f'_2(x_0),
\]
ou même :
\[
(7\text{ ter})\quad (f_1\cdot f_2)'=f'_1\cdot f_2+f_1\cdot f'_2.
\]
Ceci s'applique notamment lorsque $F_1=F_2=F=K$ (fonctions à valeurs scalaires).

f). \uline{Théorème des fonctions implicites}. Les notations étant celles de b), on suppose que $E$ et $F$ sont \uline{normés} et \uline{complets} (ce qui est toujours le cas lorsque $E$ et $F$ sont de dimension finie, cf. EVT, chap. I, § 2), que $f$ est \uline{strictement dérivable} au point $x_0$, et que $f'(x_0):E\longrightarrow F$ est un \uline{isomorphisme} d'espaces vectoriels topologiques. Soit $y_0=f(x_0)$. Il existe alors $^{(3)}$ une application $g$ d'un voisinage de $y_0$ dans $E$, strictement dérivable en $y_0$, prenant la valeur $x_0$ en ce point, et telle que $g\circ f=1_E$ (resp. $f\circ g=1_F$) au voisinage de $x_0$ (resp. de $y_0$).$^{(4)}$

g) \uline{Dérivées partielles}. On suppose maintenant que $E=K^n$ ; soit $x_0=(x_1^0,\ldots,x_n^0)\in E$. La fonction $f:U\to F$ peut donc être considérée comme une fonction $f(x_1,\ldots,x_n)$ de $n$ variables scalaires, à valeurs dans $F$, définie lorsque les $x_i$ sont assez voisins des $x_i^0$. Si on laisse les $x_j$ fixes, égaux à $x_j^0$, pour tout $j\neq i$, on obtient une fonction de la seule variable $x_i$, définie au voisinage de $x_i^0$.

\mspage{260}
